\documentclass[conference]{IEEEtran}

\usepackage{amsmath}
\usepackage{booktabs}
\usepackage{graphicx}
\usepackage{url}

\graphicspath{{../../../data_process/attnres/pic/}}

\title{Interface-Preserving Online Softmax Merge Offload for Spatz}

\author{
  \IEEEauthorblockN{Wenxuan Tang}
  \IEEEauthorblockA{Local Spatz Prototype Study}
}

\begin{document}

\maketitle

\begin{abstract}
Block-wise attention implementations rely on online softmax merge recurrences
to combine partial row summaries without materializing full score matrices.
The merge is small enough to be a poor fit for heavyweight acceleration, but
structured enough to be expensive when repeatedly executed as scalar software:
it mixes scalar maxima, scalar normalizers, vector output state, fine-grained
memory traffic, and a strict per-row update order. This paper presents a
two-stage online softmax merge offload study on Spatz. Phase A integrates a
restricted-semantics merge-update engine through memory-mapped control and a
TCDM master port, validating that the cluster can host the offload path without
changing the Spatz ISA or vector pipeline. Phase B keeps the same software
visible interface but replaces the restricted behavior with a complete
mixed-scalar merge datapath using fixed 256-segment exponential and reciprocal
lookup approximations plus fixed-point weighted accumulation. In Verilator RTL
microbenchmarks, the complete datapath passes full-reference correctness
gates, reaches up to 37.89$\times$ speedup over an RTL-aligned software
reference for the evaluated full-merge sweep, and records a worst observed
absolute error below $10^{-5}$ on the tested numeric workloads. The result is a
low-intrusion merge-offload building block for future attention acceleration
work, not a claim of end-to-end attention speedup.
\end{abstract}

\section{Introduction}

Modern exact attention kernels often use block-wise softmax formulations to
avoid materializing all intermediate scores. In these formulations, each row
maintains a compact online state consisting of a running maximum $m$, a running
normalizer $l$, and an output vector $O$. When a new tile is processed, its
partial summary is merged into the old state. This recurrence is central to
online softmax and IO-aware attention designs~\cite{online-softmax2018,
flashattention2022}.

Although the merge state is compact, the operation is not a simple vector
copy. A complete merge must compare scalar maxima, evaluate exponential scale
factors, combine scalar normalizers, form normalization weights, and apply
those weights to every output element. This shape is awkward for a compact
cluster path optimized around regular vector work: the arithmetic is
stateful, the scalar and vector parts are tightly coupled, and the memory
traffic is fine grained.

This paper asks whether online softmax merge can be offloaded through a small
cluster-local engine while preserving the existing processor interface. The
study is performed on Spatz, a compact RISC-V vector cluster platform built
around Snitch-style scalar cores and vector units~\cite{snitch2020,spatz2023}.
The design deliberately avoids modifying the ISA, decoder, controller, vector
functional units, vector register file, vector load/store unit, or instruction
pipeline. Instead, software configures a merge engine through existing-style
memory-mapped registers, and the engine reads and writes tightly coupled data
memory (TCDM) through its own master port.

The work proceeds in two stages. Phase A establishes the integration path with
a restricted merge semantics. It validates the control path, TCDM traffic,
status/error handling, software benchmark flow, and the break-even behavior of
cluster-local offload. Phase B then keeps the same interface and implements
the complete mixed-scalar online softmax merge recurrence with an approximate
datapath.

The paper makes three contributions:

\begin{itemize}
  \item It presents an interface-preserving online softmax merge offload path
        for Spatz that does not require ISA or main pipeline changes.
  \item It shows a two-stage design progression: a restricted-semantics
        integration baseline followed by a full mixed-scalar datapath
        implementing the complete online softmax merge recurrence.
  \item It provides RTL-level correctness and evaluation evidence, including
        full-reference gates, error-path regressions, cycle and TCDM counters,
        attention-like merge-chain execution, and numerical approximation
        analysis.
\end{itemize}

The scope is also explicit. The evaluation is a merge-engine study. It does
not include QK or PV matrix multiplication, external DRAM scheduling, runtime
attention orchestration, synthesis, area, power, or post-synthesis timing. The
attention-like workload is a real SMU merge-chain benchmark, but it is not an
end-to-end attention kernel.

\section{Background}

\subsection{Online Softmax Merge}

Given an old online state $(m_{old}, l_{old}, O_{old})$ and a tile state
$(m_{tile}, l_{tile}, O_{tile})$, the exact online softmax merge recurrence is:

\begin{align}
m_{new} &= \max(m_{old}, m_{tile}), \\
l_{new} &= l_{old}e^{m_{old}-m_{new}} +
           l_{tile}e^{m_{tile}-m_{new}}, \\
O_{new} &= O_{old}
          \frac{l_{old}e^{m_{old}-m_{new}}}{l_{new}} +
          O_{tile}
          \frac{l_{tile}e^{m_{tile}-m_{new}}}{l_{new}}.
\end{align}

The recurrence is structurally composable across row summaries and block
chains: a row state can be repeatedly updated by streaming additional tile
states through the same merge operator. In exact arithmetic this is the basis
for online softmax state reuse. In finite-precision hardware, this structure
still provides a useful design target, but the implementation must explicitly
control approximation error and corner-case behavior.

\subsection{Spatz Integration Context}

Spatz clusters compact RISC-V-based vector units and tightly coupled memory to
maximize computing efficiency~\cite{spatz2023}. This work uses Spatz as a
constrained validation platform. That constraint is important: the prototype
does not rely on a new vector instruction or a redesigned vector pipeline. The
merge engine is an auxiliary cluster-local unit controlled by memory-mapped
registers and connected to TCDM as a master.

This makes Spatz a minimum-integration testbed rather than a performance
ceiling. If the merge recurrence can be hosted through this conservative path,
the same datapath structure can later be widened, replicated, or moved closer
to a streaming attention front end.

\section{Design Overview}

\subsection{Software-Visible Interface}

Both phases use the same software-visible interface. Software provides base
addresses for old state, tile state, and destination state, then configures
the row count $N$, vector dimension $D$, and row stride. A control register
issues \texttt{start} and \texttt{clear\_done} pulses. A status register
reports \texttt{busy}, \texttt{done}, and \texttt{error}. The interface is
summarized in Table~\ref{tab:mmio}.

\begin{table}[t]
\centering
\caption{Memory-mapped merge engine interface.}
\label{tab:mmio}
\begin{tabular}{ll}
\toprule
Register & Purpose \\
\midrule
\texttt{MERGE\_SRC\_M\_OLD} & Base address of old $m$ array \\
\texttt{MERGE\_SRC\_L\_OLD} & Base address of old $l$ array \\
\texttt{MERGE\_SRC\_O\_OLD} & Base address of old $O$ matrix \\
\texttt{MERGE\_SRC\_M\_TILE} & Base address of tile $m$ array \\
\texttt{MERGE\_SRC\_L\_TILE} & Base address of tile $l$ array \\
\texttt{MERGE\_SRC\_O\_TILE} & Base address of tile $O$ matrix \\
\texttt{MERGE\_DST\_M} & Base address of output $m$ array \\
\texttt{MERGE\_DST\_L} & Base address of output $l$ array \\
\texttt{MERGE\_DST\_O} & Base address of output $O$ matrix \\
\texttt{MERGE\_N} & Number of rows \\
\texttt{MERGE\_D} & Vector dimension per row \\
\texttt{MERGE\_STRIDE} & Row stride in bytes; zero means packed \\
\texttt{MERGE\_CTRL} & Start and clear-done pulses \\
\texttt{MERGE\_STATUS} & Busy, done, and error bits \\
\bottomrule
\end{tabular}
\end{table}

Keeping this interface fixed across the two phases is a central design choice.
Phase A validates that the interface can drive a working cluster-local engine.
Phase B then demonstrates that the same integration path can support the full
online merge recurrence without expanding the software control contract.

\subsection{Streaming Engine Structure}

The engine is organized as a row-streaming state machine. It validates the
configuration, loads scalar state for one row, computes scalar merge values,
stores the updated scalars, streams $D$ vector elements, writes updated vector
outputs, and advances to the next row. Invalid dimensions, misaligned
addresses, misaligned strides, and unsupported numeric cases report
\texttt{MERGE\_STATUS.error} rather than silently producing results.

Figure~\ref{fig:engine-flow} shows the Phase-A control flow. Phase B retains
the same scheduling structure and replaces the restricted computation blocks
with the full mixed-scalar datapath described in
Section~\ref{sec:phase-b}.

\begin{figure}[t]
\centering
\includegraphics[width=\linewidth]{online_softmax_merge_engine_flow.png}
\caption{Streaming merge-update engine control flow used to validate the
cluster-local offload path.}
\label{fig:engine-flow}
\end{figure}

\section{Phase A: Restricted Integration Baseline}

Phase A answers a narrow but necessary question: can this kind of merge offload
be integrated safely into Spatz without modifying the ISA or the main
pipeline? The Phase-A datapath intentionally implements a restricted semantics
rather than the full online softmax recurrence. It supports zero-length old or
tile states and constructed equal-weight cases where the scalar summaries have
matching maxima and normalizers. General mixed-scalar inputs are reported as
unsupported errors.

This restriction is deliberate. It isolates the system concerns that must be
solved before a full datapath is useful: memory-mapped configuration,
completion and error signaling, TCDM request scheduling, software benchmark
integration, and performance counter observation.

\subsection{Phase-A Correctness Scope}

The Phase-A benchmark checks supported outputs against a CPU reference for the
restricted semantics. It also checks invalid configurations such as
\texttt{N=0}, \texttt{D=0}, misaligned addresses, and misaligned strides.
Legal but unsupported mixed-scalar probes are expected to report an error.

This conservative error policy prevents the restricted prototype from making
false claims about the complete online softmax equation. In the combined
article, Phase A should therefore be understood as an integration baseline,
not as a final softmax merge implementation.

\subsection{Phase-A Evidence}

The restricted engine shows the expected offload behavior: tiny cases are
slower because memory-mapped setup and polling dominate, while larger cases
amortize the fixed cost. In the representative restricted equal-weight case,
the engine reaches 2.40$\times$ speedup at $N=16,D=64$, corresponding to a
58.3\% cycle reduction. A finer break-even sweep places the threshold between
roughly 64 and 96 vector elements.

Figure~\ref{fig:a-speedup} summarizes the Phase-A speedup curve, while
Figure~\ref{fig:a-tcdm} confirms that the engine uses its TCDM master path and
that congestion remains measurable in the single-engine prototype.

\begin{figure}[t]
\centering
\includegraphics[width=\linewidth]{online_softmax_merge_bypass_speedup.png}
\caption{Phase-A restricted-semantics speedup over scalar software.}
\label{fig:a-speedup}
\end{figure}

\begin{figure}[t]
\centering
\includegraphics[width=\linewidth]{online_softmax_merge_bypass_tcdm.png}
\caption{Phase-A TCDM accessed and congested counters.}
\label{fig:a-tcdm}
\end{figure}

\section{Phase B: Full Mixed-Scalar Datapath}
\label{sec:phase-b}

Phase B answers the next question: can the same integration path support the
full mixed-scalar online softmax merge equation? The Phase-B RTL keeps the
existing \texttt{MERGE\_*} interface and the same TCDM path, but replaces the
restricted scalar/vector behavior with a complete datapath.

\subsection{Scalar Merge}

The scalar stage computes:

\begin{align}
m_{new} &= \max(m_{old}, m_{tile}), \\
s_{old} &= l_{old}e^{m_{old}-m_{new}}, \\
s_{tile} &= l_{tile}e^{m_{tile}-m_{new}}, \\
l_{new} &= s_{old} + s_{tile}.
\end{align}

The implementation accepts finite normal FP32 values and zeros within the
engine's supported numeric contract. Both-zero length inputs are treated as an
unsupported error path, because the normalized output weights would be
undefined.

\subsection{Approximation Units}

The exponential path uses a fixed Q1.23 lookup table with 256 linear segments
over $[-8,0]$. Inputs above the range saturate to one, and inputs below the
range saturate to zero. The reciprocal path normalizes the input mantissa to
$[1,2]$ and uses a 256-segment Q1.23 lookup table with interpolation. No
Newton refinement is used in the current configuration.

The datapath uses fixed-point helper functions for FP32 conversion, Q16.32
scaling, and weighted accumulation. The vector stage computes:

\begin{align}
w_{old} &= s_{old}/l_{new}, \\
w_{tile} &= s_{tile}/l_{new}, \\
O_{new} &= w_{old}O_{old} + w_{tile}O_{tile}.
\end{align}

This creates a complete mixed-scalar merge datapath while preserving the
Phase-A software interface.

\subsection{Correctness Methodology}

The software benchmark contains an RTL-aligned reference model for the fixed
approximation configuration. The full-reference probe that was an expected
error in Phase A is converted into a correctness gate in Phase B. The
benchmark also retains invalid-configuration checks, the both-zero length
error path, stride-zero packed layout coverage, and the restricted-semantics
regression cases.

The final RTL validation command was:

\begin{verbatim}
cd hw/system/spatz_cluster/sw/build
ctest -R online-softmax-merge -V
\end{verbatim}

The test passed with \texttt{online-softmax-merge PASS}. The same flow also
builds through the Verilator system binary and software targets.

\section{Evaluation}

\subsection{Full Mixed-Scalar Sweep}

Table~\ref{tab:full-mixed} reports the complete mixed-scalar sweep. The CPU
baseline is an RTL-aligned fixed-point software reference used for
same-semantics comparison and correctness gating. It is not an optimized
libm, SIMD, vector, or end-to-end attention baseline.

\begin{table}[t]
\centering
\caption{Phase-B full mixed-scalar SMU microbenchmark.}
\label{tab:full-mixed}
\begin{tabular}{rrrrr}
\toprule
$N$ & $D$ & CPU cycles & SMU cycles & Speedup \\
\midrule
1 & 1 & 2327 & 1207 & 1.93$\times$ \\
4 & 8 & 17490 & 1478 & 11.83$\times$ \\
8 & 16 & 56103 & 2343 & 23.94$\times$ \\
8 & 32 & 97554 & 3371 & 28.94$\times$ \\
16 & 64 & 365620 & 9650 & 37.89$\times$ \\
\bottomrule
\end{tabular}
\end{table}

The largest tested point reduces cycles from 365620 to 9650:

\begin{equation}
\frac{365620 - 9650}{365620} = 97.36\%.
\end{equation}

The speedup increases with workload size because the fixed launch and polling
costs are amortized over more vector elements and because the scalar software
reference performs the same approximation and fixed-point operations
explicitly.

\subsection{TCDM Counters}

Table~\ref{tab:tcdm} reports TCDM accessed and congested counters for the same
sweep. Accesses increase with workload size, as expected for a streaming
engine that reads old/tile state and writes merged state through TCDM.
Congestion is observable but small relative to accesses in this single-engine
microbenchmark.

\begin{table}[t]
\centering
\caption{Phase-B TCDM counters.}
\label{tab:tcdm}
\begin{tabular}{rrrr}
\toprule
$N$ & $D$ & Accessed & Congested \\
\midrule
1 & 1 & 327 & 0 \\
4 & 8 & 495 & 2 \\
8 & 16 & 985 & 4 \\
8 & 32 & 1583 & 9 \\
16 & 64 & 5229 & 27 \\
\bottomrule
\end{tabular}
\end{table}

\subsection{Numerical Approximation}

A host numeric model compares the fixed approximation configuration against a
full FP32 reference on representative mixed-scalar and delta-sweep workloads.
Table~\ref{tab:error} summarizes the recorded error metrics. The worst
observed maximum absolute error is $9.775\times10^{-6}$, the worst guarded
maximum relative error is $3.325\times10^{-6}$, and the worst mean absolute
error is $2.174\times10^{-6}$.

\begin{table}[t]
\centering
\caption{Numeric approximation summary.}
\label{tab:error}
\begin{tabular}{lrrr}
\toprule
Case & Max abs. & Guarded rel. & Mean abs. \\
\midrule
$4\times8$ mixed & $3.58e{-7}$ & $3.58e{-7}$ & $8.16e{-8}$ \\
$8\times16$ mixed & $2.50e{-6}$ & $2.46e{-6}$ & $4.46e{-7}$ \\
$8\times32$ mixed & $4.41e{-6}$ & $2.53e{-6}$ & $8.41e{-7}$ \\
$16\times64$ mixed & $9.78e{-6}$ & $2.58e{-6}$ & $2.17e{-6}$ \\
$16\times64$ delta & $6.08e{-6}$ & $3.33e{-6}$ & $8.10e{-7}$ \\
\bottomrule
\end{tabular}
\end{table}

These results meet the initial $10^{-3}$ target and the exploratory
$10^{-4}$ target on the evaluated inputs. They should not be interpreted as a
global error proof over all attention distributions. The input space includes
different logit deltas, length magnitudes, old/tile weight ratios, output
signs, block orders, mask patterns, sequence lengths, and multi-block error
accumulation. A broader claim would require a mathematical error bound or a
larger stress suite over realistic and adversarial traces.

\subsection{Attention-Like Merge Chain}

Phase B also includes an attention-like fallback workload that repeatedly
invokes the real SMU hardware path over multiple rows, multiple merge blocks,
and a hidden dimension. The evaluated configuration uses 8 rows, 4 blocks, and
$D=32$. It records 388826 software-reference cycles and 13432 SMU cycles,
corresponding to 28.95$\times$ speedup, with 6188 TCDM accessed events and 32
TCDM congested events.

This workload is useful because it exercises a chain of online merge state
updates. However, it is not an end-to-end attention benchmark. It excludes QK
score generation, PV accumulation, DRAM behavior, runtime scheduling, and
model-level accuracy.

\section{Scalability Discussion}

The contribution is not Spatz-only in principle, but the evidence is
Spatz-prototype evidence. The merge recurrence is structurally reusable across
rows, block chains, and heads by streaming different $(m,l,O)$ tuples through
the same datapath. The current implementation serializes vector elements
through one compact engine, but the datapath has natural replication points:
vector lanes, row-level parallel engines, independent head-level engines, and
a streaming or DMA-like front end that reduces MMIO launch overhead.

The current Spatz prototype is intentionally conservative. It uses the
existing \texttt{MERGE\_*} MMIO/TCDM interface and avoids all ISA or pipeline
changes. These constraints limit launch overhead, available engine
parallelism, and memory bandwidth in this implementation. They also
demonstrate the minimum integration requirements of the merge-offload pattern.

The measured trend supports further scaling work: engine cycles remain
favorable relative to the RTL-aligned software reference across the evaluated
sizes, and TCDM congestion is observable and measurable rather than
functionally blocking. This suggests engineering value in widening the vector
update stage, replicating engines, or replacing polling with a lower-overhead
streaming front end. It does not prove linear scaling for large end-to-end
attention workloads.

\section{Limitations}

The current study has five main limitations.

First, the evaluation is not an end-to-end attention speedup result. The
attention-like workload is a real SMU merge-chain benchmark, but it does not
compute complete attention.

Second, the CPU baseline in the Phase-B tables is an RTL-aligned fixed-point
software reference. This is the right baseline for same-semantics correctness
and microarchitectural comparison, but it is not a tuned CPU or vector
attention implementation.

Third, the numeric validation is representative rather than exhaustive. The
observed error is small on the evaluated workloads, but the paper does not
claim a proof over all softmax or attention input distributions.

Fourth, the approximation configuration is fixed at 256 exponential segments
and 256 reciprocal segments. Parameter sweeps, Newton-refinement variants, and
area/accuracy tradeoffs remain future work.

Fifth, synthesis, area, power, and post-synthesis timing data are missing.
Consequently, this paper should be read as an RTL system-prototype and
microbenchmark study rather than a final PPA comparison.

\section{Related Work}

\textbf{Online softmax.}
Milakov and Gimelshein introduced online normalizer calculation for softmax,
showing that the normalizer can be maintained in a streaming fashion while
reducing memory traffic~\cite{online-softmax2018}. This work targets the
hardware realization of the merge recurrence that appears in such online
formulations.

\textbf{IO-aware attention.}
FlashAttention restructures exact attention around tiling and IO-awareness,
reducing traffic between high-bandwidth memory and on-chip SRAM
~\cite{flashattention2022}. The present work is narrower: it does not build a
complete attention kernel, but studies a cluster-local hardware path for the
online merge state that such kernels repeatedly update.

\textbf{Compact vector clusters.}
Snitch explores tiny scalar cores for efficient floating-point workloads
~\cite{snitch2020}, and Spatz builds compact RISC-V vector units around this
cluster philosophy~\cite{spatz2023}. The SMU follows a different integration
point from ISA or vector-pipeline extension: it is an auxiliary TCDM master
controlled through memory-mapped registers.

\section{Conclusion}

This paper presented a two-stage online softmax merge offload study on Spatz.
Phase A demonstrates that the merge offload path can be integrated through a
conservative auxiliary engine interface, validating MMIO control, TCDM access,
status/error handling, and restricted-semantics benchmark flow. Phase B shows
that the same interface can support the complete mixed-scalar online softmax
merge recurrence with fixed LUT-based approximation units and fixed-point
weighted accumulation.

The final prototype passes RTL correctness gates, preserves the existing
\texttt{MERGE\_*} software interface, reaches up to 37.89$\times$ speedup over
an RTL-aligned software reference on the evaluated full mixed-scalar sweep,
and observes worst-case absolute error below $10^{-5}$ on the tested numeric
workloads. These results support the design as a low-intrusion
merge-offload building block for future attention acceleration work, while
leaving end-to-end attention integration, broader numeric validation,
parameter sweeps, and PPA evaluation as future work.

\bibliographystyle{IEEEtran}
\bibliography{../../common/bib/refs}

\end{document}

